% Bab IV  Tata Cara Penyusunan dan Ujian Skripsi
\chapter{TATA CARA PENYUSUNAN DAN UJIAN SKRIPSI}

Bab ini menguraikan alur penyusunan skripsi dari pengajuan judul sampai penjilidan. Jumlah pertemuan bimbingan, batas waktu, dan persyaratan administrasi yang bersifat teknis ditetapkan oleh Program Studi dan dapat berubah sesuai kebijakan; mahasiswa wajib menanyakan ketentuan terbaru kepada Program Studi.

\section{Alur Umum}
Tahapan penyusunan skripsi secara berurutan adalah sebagai berikut:
\begin{enumerate}
\item pengajuan judul penelitian;
\item penetapan dosen pembimbing;
\item penyusunan proposal dan bimbingan;
\item seminar proposal;
\item pelaksanaan penelitian dan bimbingan;
\item penyusunan laporan hasil dan seminar hasil;
\item penyusunan skripsi dan ujian skripsi (ujian akhir);
\item revisi, pengesahan, penjilidan, dan pengumpulan.
\end{enumerate}
Setiap tahap memiliki dokumen dengan format tersendiri, yaitu Formulir Pengajuan Judul, Template Proposal Skripsi, Format Seminar Hasil, dan Format Skripsi.

\section{Pengajuan Judul}
Mahasiswa mengajukan judul penelitian menggunakan Formulir Pengajuan Judul Penelitian yang memuat tiga usulan judul. Setiap usulan judul disertai:
\begin{enumerate}
\item judul (huruf kapital, paling banyak 15 kata);
\item objek atau tempat penelitian beserta alamatnya;
\item latar belakang masalah singkat;
\item rumusan masalah dalam bentuk penomoran;
\item tujuan penelitian yang menjawab rumusan masalah;
\item penelitian terkait minimal 10 artikel jurnal atau prosiding, lengkap dengan persamaan dan perbedaannya dengan penelitian yang diusulkan;
\item daftar pustaka yang dirujuk.
\end{enumerate}
Formulir diverifikasi oleh Ketua Program Studi atau dosen verifikator. Judul yang disetujui ditetapkan beserta dosen pembimbing I dan II.

\section{Bimbingan}
Mahasiswa wajib berkonsultasi secara berkala kepada dosen pembimbing selama penyusunan proposal, pelaksanaan penelitian, dan penulisan skripsi. Ketentuan bimbingan:
\begin{enumerate}
\item setiap bimbingan dicatat pada lembar atau kartu bimbingan yang ditandatangani dosen pembimbing;
\item naskah yang diserahkan untuk dibimbing harus sudah sesuai format penulisan pada Bab III;
\item perbaikan dari bimbingan sebelumnya dikerjakan sebelum bimbingan berikutnya;
\item perubahan judul atau ruang lingkup penelitian harus mendapat persetujuan tertulis dosen pembimbing dan Program Studi.
\end{enumerate}

\section{Seminar Proposal}
Proposal yang telah disetujui dosen pembimbing diseminarkan di hadapan dosen penguji dan mahasiswa. Naskah proposal memuat: sampul, lembar persetujuan, abstrak, daftar isi, Bab I Pendahuluan, Bab II Tinjauan Pustaka, Bab III Metode Penelitian (termasuk jadwal pelaksanaan), dan daftar pustaka. Hasil seminar berupa masukan perbaikan yang harus dituntaskan sebelum penelitian dilaksanakan.

\section{Pelaksanaan Penelitian}
Penelitian dilaksanakan sesuai proposal yang telah disetujui. Data dan catatan penelitian disimpan dengan baik karena dapat diminta oleh pembimbing atau penguji untuk diperiksa. Bila terjadi perubahan metode atau kondisi yang berpengaruh pada hasil, perubahan itu dikonsultasikan dengan dosen pembimbing dan dicantumkan dalam skripsi.

\section{Seminar Hasil}
Seminar hasil dilaksanakan setelah penelitian selesai dan Bab I sampai Bab V telah tersusun. Naskah seminar hasil memuat lembar persetujuan, abstrak, daftar isi, daftar tabel, daftar gambar, daftar lampiran, Bab I sampai Bab V, daftar pustaka, dan lampiran. Seminar hasil bertujuan memperoleh masukan atas hasil dan pembahasan sebelum ujian akhir.

\section{Ujian Skripsi}
Ujian skripsi adalah ujian lisan di hadapan dosen penguji dan pembimbing untuk menilai kemampuan mahasiswa mempertanggungjawabkan skripsinya. Penilaian meliputi isi dan mutu ilmiah skripsi, penguasaan materi dan metode, kemampuan presentasi dan menjawab pertanyaan, serta kepatuhan terhadap format penulisan. Mahasiswa dinyatakan lulus, lulus dengan perbaikan, atau belum lulus sesuai keputusan tim penguji.

\section{Revisi dan Pengesahan}
Mahasiswa yang dinyatakan lulus dengan perbaikan wajib memperbaiki skripsi sesuai masukan penguji dalam batas waktu yang ditetapkan Program Studi. Perbaikan disetujui secara tertulis oleh penguji dan pembimbing. Setelah itu skripsi disahkan dengan tanda tangan Pembimbing I, Pembimbing II, Penguji, Ketua Program Studi, dan Dekan Fakultas Teknik pada halaman pengesahan.

\section{Penjilidan dan Pengumpulan}
Skripsi yang telah disahkan dijilid dengan jilid skripsi standar berwarna biru laut yang dilaminasi. Jumlah eksemplar yang dikumpulkan, serta berkas digital (PDF) yang harus diserahkan ke perpustakaan dan Program Studi, ditetapkan oleh Program Studi.
